\section{Positions and Frontiers}
\label{sec:frontier}

The previous sections traced where visual generation \emph{is} today, level by level and pipeline by pipeline. This section turns to where the field is heading. Our organizing principle is the same five-level taxonomy of \Cref{sec:evolution}: each subsection picks an open frontier, names the unsolved problem at its center, and locates that problem within the L1--L5 progression. The intent is not exhaustive coverage but a small set of positions that, in our reading, will dominate the next two years of research on image generation and editing.

\subsection{Visual Chain-of-Thought (vCoT): Thinking Before Rendering}
\label{subsec:vcot}

The text-to-pixel mapping at the heart of every Level~1 generator is a single jump: it commits to a final image without any intermediate representation that the user, or the model itself, can inspect, debate, or revise. \emph{Visual Chain-of-Thought} (vCoT) attempts to break this jump into a sequence of intermediate states---textual analyses, layout sketches, regional decompositions, or coarse renderings---that the model produces before the final pixels are ever committed. The reason this matters mirrors why text CoT mattered for language models: complex visual outputs are easier to plan when the planning artefact is exposed to constraint checking, retrieval, or explicit revision. Several systems already realize restricted forms of vCoT---ReasonGen-R1~\citep{zhang2025reasongen} and T2I-R1~\citep{jiang2025t2i} let an autoregressive generator emit a reasoning trace before token-level synthesis, while X-Planner~\citep{Yeh2025BeyondSE} and MIRA~\citep{Zeng2025MIRAMI} translate user intent into stepwise edit sub-goals before any pixel changes. The frontier question is no longer whether reasoning helps, but how to design intermediate states that are simultaneously faithful to the eventual image and cheap enough to revise.

A first axis of design is the modality of the intermediate state. Pure-text reasoning is computationally cheap and inherits directly from the maturity of LMM-style chain-of-thought, but it is also lossy: many spatial relations and aesthetic judgements cannot be faithfully described in language without becoming verbose enough to be useless. ReasonEdit~\citep{Yin2025ReasonEditTR} and LegoEdit~\citep{Jia2025LegoEditAG} push toward typed symbolic edit programs that are richer than free-form text but still discrete; ImageEA~\citep{Hu2025ImageEA} goes further by translating natural language into executable visual operations whose effects can be verified step by step. A complementary line treats coarse visual sketches---layouts, depth proxies, or low-resolution drafts---as the intermediate representation: ReCon~\citep{ReCon} grounds compositional instructions through a visual reasoning trace, while CreatiLayout~\citep{CreatiLayout} and HybridLayout~\citep{wu2025hybrid} promote the layout itself to a first-class artefact that downstream rendering must respect, and Make Geometry Matter~\citep{zhang2026make} pushes this further by elevating explicit geometric structure as the planning substrate when 2D layouts cannot capture the spatial relations at stake. The trade-off mirrors the broader text-vs-image debate in multimodal reasoning: text traces are easier to learn and verify but information-poor; visual traces preserve geometry but are harder to score automatically.

Despite enthusiastic adoption, explicit vCoT does not always improve final image quality, and a healthy debate is emerging about when the additional latency is justified. Reasoning-augmented generation appears most useful in three regimes: (i)~when the prompt encodes hard compositional constraints---counts, spatial relations, attribute binding---that one-shot generators routinely fail, mirroring the Level-2 stress tests of \Cref{sec:stress_test}; (ii)~when the request requires external world knowledge that must be retrieved, verified, or combined, the case Gen-Searcher~\citep{feng2026gen} pushes to its logical extreme by training a search-augmented image-generation agent; and (iii)~when the user wants to inspect or steer the planning process itself, as in the agent-to-Lightroom protocol of JarvisArt~\citep{lin2025jarvisart}. Conversely, for purely aesthetic or familiar-distribution prompts, explicit reasoning often costs more than it saves and can amplify hallucinations in the trace, as the physics-exam case in \Cref{subsubsec:physics_solver_case} illustrates: the final solved image is correct, yet the trace repeatedly revisits the same sub-problems before converging.

The open challenges of vCoT therefore mirror those of text CoT but with sharper visual constraints. \emph{Faithfulness of the trace} is paramount---a sophisticated reasoning monologue that does not constrain the final pixels is worse than no reasoning, because it falsely advertises competence. \emph{Length control} matters too: many recent reasoning-augmented generators emit traces that revisit the same sub-problems repeatedly, suggesting that visual reasoning still lacks a stable internal workspace. Finally, \emph{verifiability} is largely unsolved for visual sketches: a textual reasoning step can be checked against a knowledge base, but a layout sketch can only be checked against the image it eventually produces, which makes debugging the planner difficult. Progress on vCoT will require richer intermediate representations that are simultaneously inspectable, executable, and grounded in the eventual visual output.

\subsection{Closed-Loop Visual Agents: Generation as Action}
\label{subsec:closed_loop_agents}

Across \Cref{sec:applications,subsec:agentic_visual_generation} we have seen visual generation increasingly embedded inside larger systems, but most of those systems still treat generation as an \emph{output}: the agent reasons, retrieves, plans, and only at the end commits to a final image. The next conceptual shift is to demote generation from output to \emph{action}---a primitive operation that an agent invokes mid-trajectory, observes, evaluates, and either commits or rolls back. Once generation is an action, the standard machinery of agent loops becomes available: planners decide when to generate, verifiers decide whether to keep the result, and policies learn over sequences of generate-and-revise steps rather than over single prompts. The intuition mirrors the move from one-shot LLM completion to ReAct-style tool-using assistants: the marginal gain comes not from a better generator, but from giving the model a closed loop in which to operate.

The most concrete instantiation of generation-as-action lives in the editing literature. In multi-turn editing, every operation---change the background, add an accessory, transfer a style---is an action whose precondition (current image) and post-condition (target image) are both visible to the agent. Frameworks such as X-Planner~\citep{Yeh2025BeyondSE}, MIRA~\citep{Zeng2025MIRAMI}, and ImageEA~\citep{Hu2025ImageEA} make this explicit by decomposing edit instructions into typed, executable steps. A complementary line studies \emph{transitions} themselves as a first-class action: VTG~\citep{yang2025vtg} learns to generate smooth, semantically coherent transitions between two given visual states, treating the intermediate motion as the unit of control rather than the final frame. Once transitions are addressable, an image agent can reason at the granularity of trajectories---how to morph, blend, or interpolate visual content under a high-level goal---instead of producing isolated end states.

A closed loop without verification quickly degenerates into one-shot generation in disguise. The defining capability of this family of agents is therefore not the generator, but the \emph{verifier} that decides whether an action's result is acceptable. GEMS~\citep{he2026gems} integrates a verifier and refiner inside the agent's persistent memory, allowing the system to back out of a faulty intermediate state without restarting from scratch. JarvisArt~\citep{lin2025jarvisart} exposes the same loop in a real creative application, where each Lightroom operation is reversible and inspectable before the next is committed. In embodied settings, CoT-VLA~\citep{zhao2025cotvla} and UniPi~\citep{du2023unipi} show that visual prediction can serve as a checked-out hypothesis space in which the agent simulates several plausible next states before committing to a physical action. The shared pattern is a clean separation between proposing (generation) and accepting (verification), with explicit fall-back state for the rejected path.

Even with explicit loops, three issues remain hard. \emph{Compounding error}: when an agent stitches many short generations into a long trajectory, small per-step inaccuracies accumulate, and identity drift or geometric inconsistency can become catastrophic---the same failure mode that current world models exhibit on long horizons (\Cref{subsec:world_models}) but with sharper visual signatures. \emph{Verifier reliability}: most current verifiers are themselves learned models, and a closed loop in which the verifier and the generator share a common bias is no better than open-loop generation. \emph{Latency and cost}: every additional generation in the loop multiplies inference cost, so practical systems must learn when to short-circuit a loop with a confident one-shot answer. Resolving these jointly is what separates a closed-loop visual agent from a slow open-loop generator with extra steps.

\subsection{Agentic Visual Generation: Tool-Augmented Rendering}
\label{subsec:agentic_visual_generation}

The recent rise of general-purpose agents in coding and computer-use domains points to a broader lesson for visual generation: capability is increasingly a property of the orchestration loop, not just of the underlying model. Systems such as Claude Code and Codex are powerful not because their next-token predictions are individually stronger, but because they plan, call tools, inspect external state, verify intermediate results, and iteratively refine outputs. A similar transition is now beginning in the visual domain. Instead of treating image generation as a one-shot mapping from prompt to pixels, an agentic visual system can decompose a request into retrieval, grounding, reasoning, and rendering substeps, invoking external tools whenever the answer depends on information that the base model should not be expected to memorise or hallucinate correctly.

We use the term \emph{Agentic Visual Generation} to denote this shift from passive rendering to \emph{tool-mediated visual production}. In this setting, generation is only one module inside a larger execution graph. For a single request, the agent may first query the web for real-time facts, retrieve structured data from databases, call OCR or document parsers, invoke charting or geometric engines for exact layout, use segmentation or tracking tools to localise edit regions, and finally hand the grounded plan to a visual renderer or editor. The output image is therefore not produced from latent priors alone; it is the result of a closed-loop pipeline that combines external state with learned visual priors.

Recent systems already instantiate different slices of this execution graph. GEMS~\citep{he2026gems} treats multimodal generation as an agent-native optimization loop with a planner, decomposer, verifier, refiner, persistent memory, and on-demand skill library, demonstrating that orchestration itself can unlock capability beyond the base generator. Gen-Searcher~\citep{feng2026gen} targets knowledge-intensive image requests by training a search-augmented agent that performs multi-hop web search and gathers textual evidence plus reference images before rendering, thereby grounding generation in external and potentially up-to-date state. JarvisArt~\citep{lin2025jarvisart} closes the loop on the editing side: rather than regenerating pixels from scratch, it maps multimodal user intent into executable Lightroom operations through a dedicated agent-to-software protocol, showing how visual agents can work \emph{inside} professional creative tools rather than outside them.

This tool-augmented setting is especially valuable for three failure regimes that recur throughout this work: \textbf{world knowledge}, \textbf{accuracy}, and \textbf{real-time information}. World knowledge matters when a model must generate historically correct city plans, scientifically faithful diagrams, or domain-specific dashboards. Accuracy matters when small symbolic mistakes---incorrect counts, wrong labels, misplaced overlays---destroy practical usefulness even if the image looks realistic. Real-time information matters when the requested output depends on current weather, sports standings, financial signals, software interfaces, or breaking events that cannot be reliably encoded in static model weights. In all of these cases, tool calls fundamentally change the problem formulation: instead of asking the visual model to \emph{know everything}, the agent only needs to know \emph{how to acquire, verify, and visualise} the right information.

The deeper implication is that future progress may depend less on scaling monolithic generators alone and more on building robust orchestration layers around them. The hard problem is no longer only how to render sharper pixels, but how to decide when retrieval is necessary, which tools are trustworthy, how to fuse tool outputs with visual priors, and how to verify the final image before returning it. Agentic Visual Generation is therefore not simply ``image generation with tool use.'' It is a transition from \emph{prompt-conditioned synthesis} to \emph{goal-conditioned visual task execution}. In the short term, external tools act as scaffolds that compensate for the model's missing factual, geometric, and temporal grounding; in the longer term, they may also shape the training targets for future unified systems that internalize part of this workflow. This makes Agentic Visual Generation a pragmatic bridge between today's strong but brittle renderers and tomorrow's truly grounded visual agents.

\subsection{Training with Synthetic Data and Visual Self-Play}
\label{subsec:synthetic_self_play}

A long-running thread of this work is data. \Cref{sec:resource_and_infra} documented how the field shifted from indiscriminate web scrape toward actively engineered corpora, with \Cref{sec:paradigm} tracing the rise of frontier-model distillation as the dominant data-construction recipe. Looking forward, the more provocative question is whether visual generation will increasingly close the loop and train on its own outputs---moving the field, as one early position paper put it, into the dawn of a synthetic era~\citep{yang2023aigs}. The motivation is structural rather than ideological: high-quality real images are bounded by what people choose to photograph and label, while synthetic pipelines can be steered toward exactly the failure modes the current generation of models reveals---the spatial, physical, and identity stress tests of \Cref{sec:stress_test}. The practical question is no longer whether to use synthetic data, but how to keep its distribution from collapsing under self-training.

Synthetic data in this roadmap spans three regimes whose boundaries are increasingly blurred. \emph{Frontier-model distillation} (\Cref{sec:paradigm,sec:generation}) uses a stronger generator to produce targets for a weaker one and remains the most reliable recipe when a frontier system genuinely outperforms the student. \emph{Reward-model-guided generation} pairs a generator with learned preference or correctness signals---DanceGRPO~\citep{xue2025dancegrpo}, RewardDance~\citep{wu2025rewarddance}, and EditReward~\citep{wu2025editreward} are recent examples---turning data construction into an RL-style search rather than a static teacher--student transfer. \emph{Visual self-play}, the most speculative of the three, lets the same generator both produce and critique, with retraining loops driven by self-rewarded preferences or external verifiers. Each regime trades a different axis of cost: distillation is expensive in API calls, reward-model RL is expensive in reward-model training, and self-play is expensive in the engineering required to avoid distribution collapse.

A useful complement to this high-level discussion is the concrete demand for \emph{domain-specific} synthesis where web data is sparse, culturally specific, or structurally constrained. Cross-domain transfer pipelines such as PaCaNet~\citep{PaCaNet} bridge stylistically distant Chinese painting and calligraphy through CycleGAN with transfer learning, illustrating that a small amount of paired data plus a bridging objective can manufacture training corpora for low-resource artistic domains long before diffusion-era distillation matured. More recent work in artistic character generation, exemplified by ReChar~\citep{Rechar}, focuses on \emph{structure-preserving} synthesis: each generated artefact must keep glyph integrity intact while admitting user-specified aesthetic enhancements, producing a constraint-aware synthetic-data pipeline that doubles as a benchmark for downstream training. These domain-targeted recipes are not just applications: they sketch what the next generation of synthetic-data engines must support natively---structurally constrained generation under sparse real supervision.

The frontier of this thread is visual self-play in the strict sense: a generator improved on its own filtered outputs across multiple rounds. In language modeling, similar loops have produced both impressive gains (when filtering is strict and external signals are used) and silent collapses (when the filter is itself the generator). The image-generation analogue is harder still, because visual quality is multi-dimensional and difficult to score with a single number. Promising signals come from preference-based fine-tuning---Diffusion-DPO~\citep{wallace2024diffusion}, dense-reward DPO~\citep{yang2024dense}, and video-DPO extensions~\citep{liu2025videodpo}---where a generator improves on top of its own samples filtered by human or VLM-derived preferences, but these methods still depend on an external preference signal rather than full self-play. Closed-loop self-play with no external grounding remains the open frontier of synthetic-data training, and the principal risk to manage is the one already named in the earliest synthetic-era position papers~\citep{yang2023aigs}: distribution narrowing toward whatever subset of the data manifold the generator and its self-judge already over-rate. As distillation, reward modeling, and self-play converge, the training pipeline of a frontier visual model is increasingly indistinguishable from a multi-agent system in which generators, judges, and curators continuously co-train.

\subsection{Visual Generation as World Simulation}
\label{subsec:world_models}

A recurring theme throughout this work is that visual generation is evolving from \emph{appearance synthesis} toward \emph{world modeling}: instead of producing a single static image or a passively unrolled video, the generative system must predict how a scene would evolve \emph{under intervention}. This ambition connects directly to Level~5 of our taxonomy (\Cref{subsec:level5}), where causal simulation rather than visual plausibility becomes the defining capability. We acknowledge upfront that the most vivid demonstrations in this area are video world models, which sit on the boundary of the roadmap's image-oriented scope; we include them here because they directly shape what the next generation of image generators must internalize. In this subsection, we examine the emerging frontier at which visual generators are repurposed---or purpose-built---as interactive world simulators, and discuss what this transition implies for the broader trajectory of visual intelligence.

\subsubsection{From Video Generation to Interactive Simulation}

The conceptual foundation for learning world models from visual observation dates back to the work of Ha and Schmidhuber~\citep{ha2018worldmodels}, who trained a compact latent model to predict environment dynamics and showed that policies could be learned entirely inside the model's ``dream.'' The Dreamer line of work~\citep{hafner2025dreamerv3} scaled this idea to diverse Atari and continuous-control domains, establishing that latent world models can support effective reinforcement learning with limited real-world interaction. LeCun's JEPA framework~\citep{lecun2022jepa} further articulated the theoretical case for predictive world models as the backbone of autonomous intelligence, arguing that learning to predict masked or future states in latent space is more scalable and robust than pixel-level generation.

What has changed in the past two years is the \emph{visual fidelity} and \emph{interaction bandwidth} of these models. Modern video generation backbones---DiT-based diffusion models, autoregressive transformers, and their hybrids---can now produce temporally coherent, photorealistic video at high resolution and long horizons. The natural question becomes: can these powerful visual generators be turned into \emph{playable} environments, where a user or an agent supplies actions at each step and the model renders the resulting next state in real time?

\subsubsection{Neural Game Engines and Playable Worlds}

The most vivid demonstrations of this idea come from neural game engines---systems that replace hand-coded simulators with learned generative models.

Genie~\citep{bruce2024genie} introduced a foundation world model trained on large-scale Internet video without action labels. By learning a latent action space from unlabeled video, Genie enables a user to ``play'' inside the generated environment by selecting from discovered actions, effectively creating an interactive 2D platformer from passive observation data alone. Genie~2~\citep{parker2024genie2} extended this paradigm to photorealistic 3D environments, supporting persistent memory, consistent object behavior across long horizons, and meaningful agent--world interaction at substantially higher visual fidelity than prior playable world models.

GameNGen~\citep{valevski2024gamenngen} demonstrated that a diffusion model can serve as a real-time neural game engine for the classic first-person shooter DOOM, achieving interactive frame rates while maintaining visual quality comparable to the original engine. The key insight is that once a generative model has internalized the game's dynamics, it can roll out new frames conditioned on player actions without access to any explicit game logic. DIAMOND~\citep{alonso2024diamond} pushed this direction further by showing that diffusion-based world models can match or exceed the visual fidelity of discrete-token approaches in Atari, while simultaneously serving as the environment for reinforcement learning agents.

Oasis~\citep{oasis2024} scaled the interactive world model to the domain of Minecraft, demonstrating that an autoregressive transformer can generate an open-world, first-person 3D environment in real time from keyboard and mouse inputs alone. GameGen-X~\citep{che2025gamegenx} extended interactive generation to open-world game video with support for diverse genres, multi-character scenarios, and fine-grained control over game events. Together, these systems establish a striking result: for environments with well-defined visual regularities, learned generative models can already replace hand-coded rendering and physics pipelines.

\subsubsection{World Models for Embodied and Autonomous Agents}

The same principle extends naturally to embodied settings, where visual world models serve not as entertainment engines but as simulation substrates for robot learning and autonomous driving. This line of work is closely related to the embodied applications discussed in \Cref{subsec:level5} and to the application section in \Cref{sec:applications}, but here we focus specifically on the \emph{simulation} role of world models---their capacity to stand in for the real environment and support policy optimization through imagined interaction.

In the robotics domain, UniSim~\citep{yang2024unisim} showed that a single generative model can serve as a universal simulator for visual experience, supporting both object manipulation and navigation through action-conditioned video generation. Cosmos Predict~2.5~\citep{CosmosPredict} and the broader Cosmos platform represent the industry push toward foundation-scale world models pretrained on internet-scale video and fine-tuned for embodied simulation, offering a reusable backbone for diverse downstream tasks including autonomous driving and robotic manipulation. In autonomous driving, GAIA-1~\citep{hu2024gaia1} demonstrated a generative world model that conditions on vehicle actions, text, and video to produce realistic driving scenarios, enabling controllable simulation of novel situations without real-world data collection.

A critical insight from this body of work, extensively analyzed in the companion world-model work~\citep{meietal2026videoroboticsurvey}, is that the value of a visual world model lies not in perceptual quality alone but in \emph{action faithfulness}: whether the generated future correctly reflects the causal consequences of the agent's intervention. A visually stunning rollout that ignores the commanded steering angle or misrepresents contact dynamics is worse than useless for policy learning---it actively corrupts the training signal. This requirement separates world simulation from video generation in a fundamental way: the former demands causal alignment, the latter only requires distributional plausibility.

\subsubsection{Open Challenges: From Rendering to Reasoning}

Despite the impressive demonstrations, the gap between current interactive world models and reliable world simulators remains substantial. We identify four principal challenges.

\paragraph{Long-horizon consistency.} Current models suffer from compounding errors over extended rollouts. Visual artefacts accumulate, object identities drift, and spatial layouts degrade as the generation horizon grows. Unlike traditional game engines or physics simulators that maintain explicit state, learned world models must reconstruct state implicitly from generated frames, making them vulnerable to hallucination and distributional drift.

\paragraph{Physical plausibility.} Visual world models learn correlations, not physical laws. They can approximate gravity, collision, and inertia in domains where training data is abundant, but they generalize poorly to novel physical scenarios. A model trained on platformer videos may fail on an unusual spring mechanism; a driving world model may produce implausible vehicle dynamics under extreme maneuvers. Bridging this gap likely requires integrating learned visual priors with structured physical reasoning---a direction that connects to the physics-aware generation challenges discussed in \Cref{sec:stress_test}.

\paragraph{Real-time performance.} Interactive simulation demands low-latency generation. While autoregressive models can approach real-time speeds, diffusion-based approaches typically require multiple denoising steps per frame, creating a tension between visual quality and interaction speed. The distillation and acceleration techniques discussed in \Cref{sec:method,sec:train} are directly relevant here, and further progress on few-step generation may be essential for practical deployment.

\paragraph{Compositional generalization.} Training a world model on one domain (e.g., a specific game or robot environment) does not automatically transfer to others. The most ambitious vision---a universal world simulator that can model arbitrary physical and social environments---requires compositional generalization far beyond what current models achieve. Foundation-scale approaches like Genie~2 and Cosmos represent steps in this direction, but the degree to which scale alone can solve the compositional challenge remains an open question.

\subsubsection{Implications for Visual Intelligence}

The transition from passive rendering to interactive world simulation marks a qualitative shift in what visual generation means. In a world model, an image is no longer the \emph{output}; it is a \emph{state} in a dynamic system. Generation is no longer a one-shot function call; it is one step in a closed-loop process where actions produce consequences and consequences inform future actions. This reframing dissolves the traditional boundary between ``generation'' and ``understanding'': a model that can simulate how a scene evolves under intervention necessarily possesses a form of visual understanding that goes beyond pattern matching.

From the standpoint of this roadmap's five-level taxonomy, world simulation sits at the apex because it subsumes lower levels as special cases. A world model that can faithfully simulate physical dynamics can also render controllable images (Level~2), maintain narrative consistency (Level~3), and operate in closed-loop interaction (Level~4). Whether the field will converge on monolithic world models that learn everything end-to-end, or on modular systems that compose learned visual generators with explicit physics engines and symbolic planners, remains the defining architectural question for the next stage of visual intelligence.



\subsection{Data-Centric Visual Intelligence}
\label{subsec:data_centric}

The visual generation systems discussed above are predominantly conditioned on text prompts or reference images. Conversely, a large and growing class of real-world applications requires \emph{understanding and visualizing} structured data: semi-structured tables with complex layouts, multi-page documents interleaving text, charts, and images, or analytical queries over relational databases and free-form text. As visual systems advance toward Agentic Generation (L4) and beyond, they will increasingly operate in closed-loop settings where an agent must read heterogeneous data, reason over its structure, and render faithful visual outputs such as charts, infographics, or data-driven documents for downstream consumption. This demands a capability that current generation pipelines largely lack: parsing and reasoning over structured input before generation.

Recent work from the data management community has made concrete progress on this front. \textbf{ST-Raptor}~\citep{STRaptor} introduces the Hierarchical Orthogonal Tree (HO-Tree) for question answering over semi-structured tables with merged cells, nested headers, and containment relationships, combining vision-language layout understanding with tree-based multi-step reasoning and improving answer accuracy by up to 20\% over prior methods. \textbf{MoDora}~\citep{MoDora} addresses the complementary problem at the document level, proposing a Component-Correlation Tree (CCTree) that organises pages containing interleaved text, tables, charts, and images, and supports hybrid retrieval combining LLM-guided navigation with embedding-based search, yielding accuracy gains of up to 61\%. \textbf{FDABench}~\citep{FDABench} scales to the multi-source setting, benchmarking data agents that must reason jointly over images, relational databases, and free-form text to answer a single analytical query, underscoring the difficulty of coordinating heterogeneous modalities within a unified pipeline.

A parallel line of work directly bridges structured-data understanding and visualization synthesis. \textbf{NL2VIS}~\citep{NL2VIS} formalizes the translation of natural-language queries over structured data into chart and graph specifications, connecting the well-studied NL2SQL pipeline to visual output. \textbf{DataVisT5}~\citep{DataVisT5} further trains a unified model that jointly understands text and data visualizations, supporting bidirectional tasks such as chart question answering and natural-language-to-visualization generation within a single architecture. These efforts echo the unified understanding-and-generation trend discussed throughout this work and demonstrate that faithful visualization demands deep comprehension of both data semantics and visual encoding conventions.

Together, these works point toward a data-centric vision of visual intelligence in which generation is grounded in real-world structured data rather than curated prompts alone. Tighter integration between data management and visual generation remains a largely unexplored yet promising frontier, particularly as visual agents (\Cref{subsec:closed_loop_agents,subsec:agentic_visual_generation}) begin to treat databases, spreadsheets, and document corpora as first-class inputs to the rendering loop.

\subsection{Evaluating New Generative Capabilities}
\label{subsec:eval_frontier}

A practical consequence of the new capabilities discussed above is that current evaluation protocols are no longer adequate. Frontier systems such as Nano Banana and GPT-Image now succeed routinely on tasks that earlier benchmarks did not anticipate---rendering scientifically faithful flow diagrams, drawing chemical structures with correct bond geometry, executing dense multi-line typography in multiple scripts, and producing tutor-style annotated worked solutions on top of an existing image (the physics-exam case in \Cref{subsubsec:physics_solver_case}). FID, CLIP-Score, and isolated prompt-following metrics tell us nothing about whether the molecular structure is correct, whether the flowchart's edges respect the underlying graph, or whether the rendered Chinese characters are stroke-faithful. \Cref{sec:resource_and_infra} already documents the move from heuristic perceptual scores to VLM-as-a-judge protocols; the next step is to define benchmarks that directly probe \emph{structural correctness} for the capability classes that matter.

Three concrete directions are emerging. First, \emph{symbolic-graph correctness} for diagrammatic outputs: a generated flowchart should be evaluated against the directed graph it claims to represent (nodes, edges, labels, decision branches), not against a perceptual similarity to a reference rendering. Second, \emph{domain-grounded factuality} for scientific and technical imagery: a chemistry diagram must be parsed back into a SMILES string and matched against the prompt's ground truth; an architectural floor plan must satisfy the geometric constraints implied by the requested layout. Third, \emph{glyph-level rendering audits} for typography: rendered text should be OCR'd and exact-matched against the prompt, with edit-distance scoring for partial credit and stroke-level analysis for non-Latin scripts. Each of these directions implies a different evaluator stack---a graph parser, a symbolic chemistry validator, an OCR + tokenizer pipeline---rather than a single similarity metric. Building these evaluators is itself a research agenda, and we expect it to be the bottleneck that forces visual benchmarks to mature in step with the capabilities of \Cref{sec:method,sec:applications}.

This evaluation gap is not isolated to image generation. World-simulation models (\Cref{subsec:world_models}) face the same problem on the temporal axis: how does one verify that a rolled-out physical trajectory is causally faithful, not merely visually plausible? The convergence point is that visual benchmarks of the next generation will need to look more like compilers and theorem provers than like image-similarity metrics---a shift that mirrors how NLP evaluation moved from BLEU toward execution-based and grounded metrics. Until that shift completes, headline progress on FID and CLIP-Score will continue to overstate the field's true competence on the higher rungs of \Cref{tab:5levels}.
